\documentclass[11pt]{article}

\usepackage[margin=1in]{geometry}
\usepackage{amsmath,amssymb}
\usepackage{enumitem}
\usepackage[hidelinks]{hyperref}
\usepackage{xcolor}

\setlist[itemize]{leftmargin=1.4em}
\setlist[enumerate]{leftmargin=1.6em}

\title{SYSC 5108 Exam Study\\Assignment 2, Question 1}
\author{Jesse Levine}
\date{}

\begin{document}
\maketitle

\section*{Question}

Assume that the prior for \(q\), the proportion of spam emails sent, is a random
variable with a uniform distribution. Let \(x\) denote the number of spam emails
received in a random sample of \(n\) emails. Find the posterior distribution
\(p(q \mid x)\). Which distribution will \(p(q \mid x)\) be?

\section*{Course and Textbook Cross-References}

\begin{itemize}
    \item \textbf{Chapter 3 slides, Probability Review:} use Bayes' formula,
    \(p(X \mid Y) = p(X)p(Y \mid X)/p(Y)\), with the denominator acting as a
    normalizing constant.
    \item \textbf{Chapter 3 slides, Maximum Likelihood Estimator:} the likelihood
    is the probability of the observed sample as a function of the unknown
    parameter. Here that parameter is the spam probability \(q\).
    \item \textbf{Goodfellow, Bengio, and Courville, \emph{Deep Learning},
    Chapter 3:} a Bernoulli random variable is binary and has
    \(P(X=1)=\phi\), \(P(X=0)=1-\phi\). A spam/not-spam email is therefore a
    Bernoulli trial with parameter \(q\), and \(x\) spam emails in \(n\) trials
    gives a binomial likelihood.
    \item \textbf{Goodfellow et al., Chapter 5, Bayesian Statistics:} before
    observing data, knowledge about an unknown parameter is represented by a
    prior \(p(\theta)\); after observing data, the likelihood is combined with
    the prior using Bayes' rule to obtain the posterior \(p(\theta \mid \text{data})\).
\end{itemize}

\section*{Step-by-Step Solution}

\subsection*{1. Identify the random variables}

Each email has two possible outcomes:
\[
    \text{spam} = 1, \qquad \text{not spam} = 0.
\]
If \(q\) is the unknown probability that an email is spam, then each email can be
modeled as a Bernoulli trial:
\[
    Y_i \mid q \sim \operatorname{Bernoulli}(q).
\]
The observed statistic is the count of spam emails:
\[
    x = \sum_{i=1}^{n} Y_i.
\]
Therefore, conditional on \(q\), the count \(x\) follows a binomial distribution:
\[
    X \mid q \sim \operatorname{Binomial}(n,q).
\]

\subsection*{2. Write the likelihood}

The binomial likelihood is
\[
    p(x \mid q)
    = \binom{n}{x} q^x(1-q)^{n-x},
    \qquad 0 \le q \le 1.
\]
When treating this as a function of \(q\), the binomial coefficient
\(\binom{n}{x}\) is constant because \(n\) and \(x\) are already observed.

\subsection*{3. Write the prior}

The question states that \(q\) has a uniform prior on \([0,1]\):
\[
    q \sim \operatorname{Uniform}(0,1).
\]
This has density
\[
    p(q) = 1, \qquad 0 \le q \le 1.
\]
Equivalently,
\[
    \operatorname{Uniform}(0,1) = \operatorname{Beta}(1,1).
\]

\subsection*{4. Apply Bayes' rule}

Bayes' rule gives
\[
    p(q \mid x)
    =
    \frac{p(x \mid q)p(q)}{p(x)}.
\]
The denominator \(p(x)\) does not depend on \(q\), so it only normalizes the
posterior. Using proportionality,
\[
    p(q \mid x)
    \propto
    p(x \mid q)p(q).
\]
Substitute the likelihood and prior:
\[
    p(q \mid x)
    \propto
    \binom{n}{x} q^x(1-q)^{n-x} \cdot 1.
\]
Remove the constant \(\binom{n}{x}\), since it does not depend on \(q\):
\[
    p(q \mid x)
    \propto
    q^x(1-q)^{n-x}.
\]

\subsection*{5. Match the kernel to a Beta distribution}

A Beta distribution with parameters \(\alpha\) and \(\beta\) has density
\[
    \operatorname{Beta}(q;\alpha,\beta)
    =
    \frac{1}{B(\alpha,\beta)}
    q^{\alpha-1}(1-q)^{\beta-1},
    \qquad 0 \le q \le 1,
\]
where
\[
    B(\alpha,\beta)
    =
    \int_0^1 q^{\alpha-1}(1-q)^{\beta-1}\,dq
\]
is the Beta function.

Compare the posterior kernel
\[
    q^x(1-q)^{n-x}
\]
with the Beta kernel
\[
    q^{\alpha-1}(1-q)^{\beta-1}.
\]
Matching exponents gives
\[
    \alpha - 1 = x,
    \qquad
    \beta - 1 = n-x.
\]
Therefore,
\[
    \alpha = x+1,
    \qquad
    \beta = n-x+1.
\]

\section*{Final Answer}

\[
    \boxed{
    p(q \mid x)
    =
    \operatorname{Beta}(q; x+1, n-x+1)
    }
\]
or, written as a normalized density,
\[
    \boxed{
    p(q \mid x)
    =
    \frac{1}{B(x+1,n-x+1)}
    q^x(1-q)^{n-x},
    \qquad 0 \le q \le 1.
    }
\]

The posterior is a \textbf{Beta distribution}. In words, the uniform prior
\(\operatorname{Beta}(1,1)\) contributes one pseudo-count to each class, so after
observing \(x\) spam emails and \(n-x\) non-spam emails, the posterior becomes
\[
    q \mid x \sim \operatorname{Beta}(1+x,\;1+n-x).
\]

\section*{Exam Notes}

\begin{itemize}
    \item This is the standard Beta-binomial conjugacy result:
    \[
        \operatorname{Beta}(\alpha,\beta)
        + \operatorname{Binomial}(n,q)
        \rightarrow
        \operatorname{Beta}(\alpha+x,\beta+n-x).
    \]
    \item With the assignment's uniform prior, \(\alpha=\beta=1\), so the result
    reduces to \(\operatorname{Beta}(x+1,n-x+1)\).
    \item The posterior mean is
    \[
        \mathbb{E}[q \mid x]
        =
        \frac{x+1}{n+2}.
    \]
    This is not required to answer the question, but it is a useful exam check:
    it is a smoothed estimate of the spam proportion, slightly pulled away from
    exactly \(0\) or \(1\) by the uniform prior.
\end{itemize}

\end{document}
